\documentclass{article}

\usepackage[T1]{fontenc}

\usepackage{booktabs,graphicx,siunitx}

\begin{document}

\section{Cronbach's alpha}

\begin{tabular}{lllll}
\toprule
$\omega$ & $\alpha$ & 95\% CI & N items & N cases \\
\midrule
.63 & .63 & [0.56, 0.70] & 3 & 301 \\
\bottomrule
\end{tabular}

\begin{tabular}{lll}
\toprule
Item & Corrected item-total r & $\alpha$ if item dropped \\
\midrule
x1 & 0.45 & .51 \\
x2 & 0.37 & .61 \\
x3 & 0.48 & .46 \\
\bottomrule
\end{tabular}

\includegraphics{figures/01_cronbachs-alpha/figure_item-total-correlation.png}

$\omega$ (McDonald's omega) is the headline reliability coefficient; it assumes only a congeneric (one-factor) model and is preferred over $\alpha$ when item loadings differ (McNeish, 2018). $\alpha$ (Cronbach's alpha) assumes tau-equivalence (equal loadings) and is a lower bound when loadings differ; it is retained as a secondary/legacy coefficient because reviewers still request it. Both range 0--1; $\ge$ .70 is commonly considered acceptable, $\ge$ .80 good, $>$ .95 suggests item redundancy --- treat these as heuristics, not pass/fail thresholds. The "$\alpha$ if item dropped" column flags items that, if removed, would raise $\alpha$ --- candidates for review. A high coefficient does not prove unidimensionality.

Internal consistency was below the conventional .70 threshold, $\omega$=.63 (95\% CI [0.56, 0.70]); Cronbach's $\alpha$=.63.

Statistical basis: Cronbach's alpha (tau-equivalent reliability) (Revelle W (2024). psych: Procedures for Psychological, Psychometric, and Personality Research. Northwestern University, Evanston.); McDonald's omega as the headline reliability coefficient (preferred over alpha when loadings differ) (McNeish, D. (2018). "Thanks coefficient alpha, we'll take it from here." Psychological Methods, 23(3), 412-433.).

\end{document}